\documentclass[sigconf]{acmart}
\settopmatter{printacmref=false}
\renewcommand\footnotetextcopyrightpermission[1]{}
\acmConference[Area Seminar SS 2026]{}{July 2026}{Rostock, Germany}
\usepackage{lipsum}
\usepackage{etoc}

\title{ODE Document -- Oil Slick Detection in Sentinel-1 Imagery}

\author{Malte Grube}
\affiliation{
  \institution{University of Rostock}
  \city{Rostock}
  \country{Germany}
}
\email{malte.grube@uni-rostock.de}

\author{Suma Gonuguntla}
\affiliation{
  \institution{University of Rostock}
  \city{Rostock}
  \country{Germany}
}
\email{suma.gonuguntla@uni-rostock.de}

\author{Vishnu Kaivalya Gudapati}
\affiliation{
  \institution{University of Rostock}
  \city{Rostock}
  \country{Germany}
}
\email{vishnu.gudapati@uni-rostock.de}

\begin{document}

\maketitle

\renewcommand{\contentsname}{\textbf{Table of Contents}}
\localtableofcontents \vspace{1em}
\textbf{Guide information:} This document is based on the ODE guide proposed in the paper referenced in the footer below\footnote{\url{https://www.sciencedirect.com/science/article/pii/S1364815226000599}, accessed 21/04/2026}.

\section{Overview}

\subsection{Purpose \& Problem Definition}
\subsubsection{What is the main objective of the study (prediction, classification, explanation, etc.)?} Binary classification of oil slicks in oceans using Sentinel-1 synthetic aperature radar (SAR) imagery.
\subsubsection{What are the research questions or hypotheses addressed?} How well can deep learning detect oil slicks in Sentinel-1 SAR imagery, and how does performance change between an in-distribution random split and a geographic out-of-distribution split (Mediterranean)? An optional scope is whether a Geospatial Foundation Model (GFM) can improve over a Convolutional Neural Network (CNN) baseline, and what Explainable AI (XAI) methods reveal about the features which drive the predictions.
\subsubsection{Why is ML an appropriate approach here (what alternative / statistical baselines would miss)?} First of all, oil slicks cause severe ecological damage. Furthermore, it is expensive to monitor the oceans manually. ML-based predictions are useful for scalable detection and monitoring in this scenario. Since baseline models (CNNs) do not capture the full amount of information in the SAR images, a transformer-based chip represenation by a GFM model might bring improvements by its pretraining on geospatial data. In contrast, traditional ML methods (e.g., SVMs, Random Forests) would require handcrafted features, which might not capture the full information in the SAR images. Convolutional or even transformer-based architectures can learn spatial features directly from the SAR images, which is crucial for this task.
\subsubsection{What is the model used for and what assumptions guide this choice?} Primarily for prediction: classification of oil slick presence or absence in Sentinel-1 SAR imagery. Since SAR imagery works at night and through clouds, it is suitable for the desired task. Moreover, oil films dampen short surface waves, visible as dark patches in SAR backscatter.

\subsection{Phenomenon \& Targets}
\subsubsection{If applicable, what is the target variable?} A binary classification label, indicating the presence (1) or absence (0) of an oil slick in a given Sentinel-1 SAR image chip.
\subsubsection{What is the temporal and spatial scope?} The spatial scope includes the Mediterranean Sea as an out-of-distribution test set in the geographic split. The temporal scope is not explicitly defined, but the dataset includes SAR imagery from 2015 to 2020.
\subsubsection{What assumptions or known constraints shape the study?} SAR imagery is robust to weather and lighting conditions. The key challenge is distinguishing oil slicks from ''lookalikes'', such as natural slicks, low-wind areas, or ship wakes. Moreover, the initial dataset without any filtering is imbalanced, with more positive than negative samples. It also contains some corrupted or nodata-heavy samples, which makes the dataset even smaller after filtering, but still keeps the imbalance.

\subsection{Workflow Overview}
\subsubsection{What are the main steps from raw data to results (high-level summary)?} The workflow looks as follows: (1) Data analysis and filtering to remove invalid samples, (2) Preprocessing using z-score normalization, resizing, and data augmentation (for the training dataloader), (3) Model training: a small CNN baseline trained from scratch, and fine-tuning of the TerraMind GFM with a modified classification head (using seeding, logging, and different hyperparameter configurations), (4) Evaluation using a variety of metrics, and (5) Performing failure case analysis (confusion matrices, lookalike analysis).
\subsubsection{Which software, versions, and computational environment are used?} The project is based on Python3.12.4 (using a virtual environment using e.g. pyenv or venv). The primary deep learning framework is PyTorch, coupled with torchvision. For the GFM Terramind, the library TerraTorch is used. Data handling includes numpy, pandas, and tifffile. Metrics calculation and further evaluation are performed using scikit-learn. Visualization is done with matplotlib and seaborn, plus folium for interactive maps. The data downloading script requires the Hugging Face Hub library. Packages are listed in the \texttt{requirements.txt} file.
\subsubsection{How is reproducibility ensured (random seeds, version control)?}
Version control is maintained through GitHub\footnote{\url{https://github.com/mgrube753/AS-SS2026-Oil-Slick-Detection.git}, accessed 11/05/2026}. Seeding is applied for all random processes (excluding data splitting, which is predefined).

\section{Data} \label{sec:data}

\subsection{Source data \& Accessibility}
\subsubsection{What are the data types and their sources?} The source of the data is the OilSlick subset of the WaterBench dataset, and it consists of images and tabular metadata. The imagery is based on 2-band GeoTIFFs from Sentinel-1 SAR, with VV and VH backscatter channels. Metadata is provided in CSV and JSON formats, which are primarily captured using the Sentinel-2-based dataset (optical imagery instead of radar-based). In contrast to Sentinel-1 chips, Sentinel-2 imagery is dependent on cloud-free conditions, so the metadata includes information about cloud coverage and reflecance. Sentinel-1 chips instead are robust to weather and lighting conditions (day and night), and are used for this project task by using its 2 backscatter band channels, whereas Sentinel-2 offers 12 band channels in this dataset.
\subsubsection{Are the source data (i.e., raw data) used for the study openly accessible and / or published? If so, where?} Yes, openly accessible on Hugging Face\footnote{\url{https://huggingface.co/datasets/ayushprd/WaterBench}, accessed 01/05/2026}. To obtain the full OilSlick dataset (since a subset of data is missing on Hugging Face), users should visit the Google Drive link\footnote{\url{https://drive.google.com/file/d/1yv23B9PtfBD10j2VcYs5JCcdHOOXcI4S/view?usp=sharing}, accessed 10/06/2026} after using the data downloading script provided. Then, the imagery has to be added manually to the respective directory.
\subsubsection{What licenses, permissions, and ethical considerations apply?} The dataset is licensed under CC-BY-4.0\footnote{\url{https://creativecommons.org/licenses/by/4.0/deed.en}, accessed 10/06/2026}. Permissions and ethical considerations include the free use, distribution and adaptation of the dataset, with proper attribution to the original creator. The rights of the original creator need to be respected, additional permissions might be required, and any modifications must be made transparent.
\subsubsection{How representative are the data (coverage, gaps) and what known biases may affect inference?} The unfiltered dataset is not evenly distributed, so does the filtered one as well. Positive samples appear more frequently than negative ones in both cases. This can introduce a bias where the model might over-predict the positive class, leading to more false positives. Overall, lookalike conditions, such as low-wind areas or natural surface films, might imitate oil slicks and could present a bias towards the model's predictions. These factors may affect how well the model generalizes to unseen data. A bias for the geographic split type might also be present, since the Mediterranean is specifically chosen as an out-of-distribution test set by default.

\subsection{Structure \& Quality}
\subsubsection{How large is the dataset (records, classes)?} The unfiltered dataset contains 1363 samples (the filtered shows 1196 samples), based on 2 classes (oil slick presence vs. absence). The metadata (primarily from Sentinel-2) does not provide features for Sentinel-1 imagery, which could be used for filtering. Useful variables are \texttt{center\_lon}, \texttt{center\_lat}, \texttt{subcategory}, \texttt{label}, and path columns. Several features are not useful for the Sentinel-1 dataset, such as \texttt{nodata\_fraction}, so it is needed to calculate this feature for each S1 TIFF file. This is crucial for filtering out samples.
\subsubsection{What known data issues exist, and what quality control checks were applied?} Several TIFF files for Sentinel-1 contain large nodata areas or are corrupted, shown by zero standard deviation in their pixel values. Nodata fractions were recalculated, and samples with more than 40\% nodata or invalid pixel statistics were removed. This filtering step helps to reduce the amount of invalid or corrupted samples, so that mostly valid and informative images are used for modeling overall. For more information according to data quality issues, see Chapter \ref{sec:limit}.

\subsection{Exploratory Description}
\subsubsection{What descriptive summaries (distributions, statistics) were made to understand the data before modeling?} The data analysis notebook includes plots showing the distribution of positive and negative labels, subcategory counts, and a world map of sample locations. Histograms on dB value ranges and boxplots for VV and VH channels help visualize SAR backscatter ranges. Summary statistics such as mean, standard deviation, and percentiles were also computed to understand the data before modeling. Moreover, distribution plots for nodata fractions and its filtering process are included, so that statistics can be compared.
\subsubsection{Is an Exploratory Data Analysis (EDA) available? If so, where?} The mentioned EDA notebook is available on GitHub\footnote{\url{https://github.com/mgrube753/AS-SS2026-Oil-Slick-Detection/tree/main/20_data_analysis}, accessed 13/05/2026}.

\subsection{Splitting (anti-leakage)}
\subsubsection{How were the training, validation, and test splits defined?} Two splits are defined: (1) a random split for in-distribution data, and (2) a geographic split with the Mediterranean as an out-of-distribution set.
\subsubsection{How was overlap or ''leakage'' between these sets avoided?} By using the \texttt{splits/random/} and \texttt{splits/geographic/} folders, containing pre-defined split TXT files to avoid any overlap or leakage.
\subsubsection{What share of the data went into each set?} Before\,/\,after filtering, the random split is structured as follows: 900/793 Train, 150/127 Validation, 300/266 Test data. Moreover, the geographic split is structured slightly differently: 864/749 Train, 217/189 Validation, 150/141 Test.
\subsection{Preprocessing \& Features}
\subsubsection{What preprocessing steps were applied, and where were they fitted?} First, the data filtering was done using the EDA notebook. Afterwards, each Sentinel-1 chip is normalized using the split-specific (random/geographic) means and standard deviations of the training set. By this, the z-score normalization can be calculated for each VV and VH backscatter channel in dB. This is done after splitting to prevent data leakage. Each channel is scaled to 224x224 pixels (before calculating the z-score for a speed-up), and the two channels are stacked to create a 2-channel input tensor for both models.
\subsubsection{What features were engineered and why?} No handcrafted features were created. The models operate directly on the SAR image bands, and do feature extraction while training, validation and testing. Besides calculating the nodata fraction, only normalization values (mean and standard deviation) were computed from the training sets to standardize the input images, as described above.
\subsubsection{How were collinearity and redundancy checked?} Overall, this question might not be applicable for this project task. Later on, failure cases are analyzed instead to capture false positives and false negatives for the final set of models.
\subsubsection{Was dimensionality reduction used, and if so, how and when?} No, since CNNs and GFMs can learn spatial features directly from the SAR images. Only resizing to 224x224 pixels is applied, so that the input dimensions are suitable for both models.

\section{Execution}

\subsection{Model Choice}
\subsubsection{Which candidate algorithms or model families were considered, and why?} Baselines and GFMs were provided by the project supervisor. On the one hand, standard image-classification CNNs (ResNet-18 and 50) were proposed for OilSlick as possible baselines for training from scratch. On the other hand, TerraMind, a pre-trained GFM, is provided to assess any improvements against the baseline model. This is proposed by applying transfer learning on OilSlick, solely training a classification head which gets fed with the GFM's feature embeddings.
\subsubsection{Which final models were selected, and on what grounds (performance, resource efficiency, useful comparison)?} As the baseline, none of the ResNet models were selected. Instead, a small baseline with 4 layers was crafted to train it from scratch using 2-channel SAR input. Such a vast model would require more training data, which is not available in this project. The small CNN baseline is selected to fit the limited dataset size, and to reduce the risk of overfitting. According to GFMs, TerraMind as \texttt{terramind\_v1\_small} is selected to discover whether the its pre-training on geospatial data boosts performance for this specific task, using its \texttt{S1GRD} modality.

\subsection{Training Protocol}
\subsubsection{What objective function(s) and weighting schemes were used?} The binary cross-entropy with logits loss was used, since the classification heads end with a single-neuron sigmoid. Additionally, class weighting is crucial since the dataset is imbalanced (more positive samples included). In this case, the positive class weight for this loss function was calculated as the ratio of negative to positive samples in the training set. This way, the model is penalized more for misclassifying the minority class, by introducing a positive class weight of $<1$.
\subsubsection{How were hyperparameters tuned (grid search, learning rate scheduling)?} Grid search is applied to find the best hyperparameter configuration for both models and both splits in our project, resulting in 4 final models. The hyperparameters include 3 model-specific learning rates, and 3 model-specific weight decays (9 configurations per model). On top of this, two learning rate schedulers are applied sequentially for each training run: a linear warmup phase (3 epochs for CNN, 2 epochs for GFM) to start from 10\% and reach the desired learning rate quickly, followed by a cosine annealing scheduler for the remaining epochs to lower the learning rate to 1\% of the desired learning rate slowly. Since there are two pre-defined data splits, a total of 36 training runs (2 models $\times$ 2 splits $\times$ 9 configurations) are performed. Moreover, the pipeline has a fixed batch size of 16, 32 epochs for the CNN baseline from scrach, and 16 epochs for the pre-trained GFM, which requires a shorter but more aggressive training of the classification head only. This also includes an array of higher learning rates and lower weight decays for the GFM, in contrast to the CNN baseline.
\subsubsection{What measures were taken to stop the model from overfitting (e.g., regularization penalties, early stopping during training)?} A dropout layer for each modified classification head, and early stopping based on the validation loss with a patience of 4 epochs. L2 regularization is also applied by using AdamW as the optimizer of choice. Furthermore, to create more training data, data augmentation is applied by using random horizontal and vertical flipping. Moreover, random rotation by 90-degree steps is also considered.

\subsection{Evaluation \& Baselines}
\subsubsection{What key results have been obtained?} % todo (maybe via discussion)
\subsubsection{Which metrics were chosen (primary and secondary), and why?} Accuracy, precision, recall, F1 score, and AUROC are computed during training and evaluation. To evaluate the binary classification performance robustly across both split types, the F1 score is chosen as the primary metric, since it balances precision and recall. These 3 metrics are included in the final poster, while the secondary metrics accuracy (due to class imbalance) and AUROC are logged during training and evaluation, but omitted in the poster to keep it concise.
\subsubsection{How do validation / test results support claims?} % todo (maybe via discussion to understand the question's intent)
\subsubsection{What steps were taken to check overfitting / underfitting?} To detect overfitting / underfitting, e.g.\,training and validation loss curves are plotted and logged using mlflow during training to check for any divergence (overfitting) or convergence (with high loss: underfitting). The mentioned early stopping procedure is implemented to stop training when validation loss does not improve.

\subsubsection{How was uncertainty quantified or assessed?} By using the two given data splits (random and geographic), we can assess the models' generalization performance and robustness to these different distributions. Additionally, by using two different models (CNN baseline and GFM), we can compare their performance, regarding whether the GFM's potential benefits from pre-training on geospatial data.

\subsection{Interpretation \& Limits}
\subsubsection{How were results interpreted?} For both models, the baseline CNN and the TerraMind GFM, no XAI method is applied, due to time constraints. For future work, it would be interesting to apply XAI methods to both model architectures, e.g., Grad-CAM for the CNN or attention map analysis for the GFM. % todo: mention on top the metric comparison and failure case inspection. watch in poster table: which model performs better according to either precision, recall, or F1 score? with binary small dataset, the GFM did not improve over the CNN baseline. let's see where we add this. let's make sure that we do not say the same in 3 different subsections.

\subsubsection{Do findings align with domain knowledge, and were plausibility checks performed (e.g., expert review, comparison to external knowledge)?} During the last meeting with the project supervisor, the results were discussed and compared to the supervisor's expectations. The results align with the domain knowledge, but what is surprising is that the geographic split does not show a significant drop in performance, which was expected instead. This might be due to the fact that the final models (after grid search) are forced to do binary classification on the small OilSlick Sentinel-1 dataset. Furthermore, the GFM does not show a significant improvement over the CNN baseline, which might be due to the same reasons: limited dataset size, or the feature embeddings of the GFM are not suitable for solely a two-class classification task, to improve over a small CNN baseline.
\subsubsection{What are the known limitations, failure modes, generalization caveats, and conditions for reliable use?} \label{sec:limit} 
(Failure analysis on the test of for the 4 best models has shown that certain OilSlick images still do not meet the desired quality to be classified correctly, even after filtering.) % todo: add more here (like small dataset size, class imbalances, lookalikes)
\subsubsection{What ethical or societal issues arise from this study?} ... % todo

\subsection{FAIR Data, Code and Models}
\subsubsection{How are the produced data, code, and models made findable?} The code and checkpoints of the trained models are available on the proposed GitHub repository, including documentation for each Python file. 
\subsubsection{How are they made Accessible? (Trusted repositories, terms of use)} Via the GitHub repository, which is openly accessible and are licensed under CC-BY-4.0 as the dataset itself.
\subsubsection{How are they made Interoperable? (Community standards, open formats)} Python and Jupyter notebooks are the core of the project. PyTorch model training checkpoints in the .pt format are saved for reuse. Results are saved as CSV files to enable easy analysis.
\subsubsection{How are they made reusable? (Documentation, licensing)}  Documentation and code comments are provided in the GitHub repository to ensure proper understanding and reuse of the code. Since the OilSlick dataset is used for the project, the same CC-BY-4.0 license applies to the code and models.
\subsubsection{How is the source of the model documented (history, reuse of prior models, data sources)?} The baseline model is solely hand-crafted and trained from scratch on the provided dataset using torch\,/\,torchvision. Terramind is a pre-trained GFM loaded from the TerraTorch library, and fine-tuned on the same dataset. The Sentinel-1 SAR data subset from OilSlick is sourced from the WaterBench dataset, which is openly available on Hugging Face, as stated in Section \ref{sec:data}.
\subsubsection{How are all contributors acknowledged or credited (e.g., CRediT taxonomy)?} The contributors are listed as authors in the ODE document. The contributions are distributed as follows:\subsection*{Malte}
\begin{itemize}
    \item Contributed to the exploratory data analysis notebook (see \texttt{oilslick\_eda.ipynb}), primarily with Vishnu.
    \item Designed the initial project structure and coordinated task management across the project timeline.
    \item Developed \texttt{dataset.py}, and contributed to \texttt{dataloader.py} with Suma together.
    \item Implemented z-score normalization and percentile clipping in the notebook and dataloader script.
    \item Implemented filtered metadata and normalization statistics generation, integrated the Baseline CNN and TerraMind models into the training pipeline (includin train and runner scripts) to ensure proper usage, and added MLflow logging.
\end{itemize}
\subsection*{Suma}
\begin{itemize}
    \item Documented the \texttt{30\_experiments/} modules, covering the training pipeline, execution, evaluation, and failure analysis.
    \item Designed the scientific poster, refined the layout and technical content to comply academic presentation standards.
    \item Implemented preprocessing and augmentation transforms in \texttt{dataloader.py} for all data splits.
    \item Contributed to the implementation of filtered metadata and normalization statistics with Malte.
    \item Involved in discussions on the experimental workflow, documentation, and preprocessing pipeline.
\end{itemize}
\subsection*{Vishnu}
\begin{itemize}
    \item Contributed to exploratory data analysis with Malte, primarily focusing on value range analysis and band statistics.
    \item Implemented $90^\circ$ rotation augmentation in the dataloader.
    \item Developed the Baseline CNN, implemented weight initialization, and developed \texttt{eval.py} and \texttt{failure\_analysis.py}.
    \item Contributed to saving of filtered metadata and normalization statistics by fixing a notebook issue of the group.
    \item Implemented early stopping, checkpointing, and MLflow-based model selection for final evaluation.
\end{itemize}
\subsection*{Joint Contributions}
\begin{itemize}
    \item Designed the experimental directory structure.
    \item Collaboratively developed the ODE, focusing on data, model selection, and workflow.
    \item Discussed dataset preparation and limitations of the WaterBench subset and Hugging Face release.
\end{itemize}
\subsubsection{Who are the intended audiences, and how was documentation adapted to them?} The primary audience are geospatial AI / ML researchers. To ensure both reproducibility and clarity, the documentation shows explanations of preprocessing, the baseline architecture and modifications regarding the GFM, the training protocols, and evaluation metrics.
\end{document}
